% GENERATED FILE. Edit canonical lessons, then run npm run generate:books.
% canonical-source-hash: fnv1a64:5796203c29a8a6ce
% canonical-lessons: GE-C510-tribuene, GE-C510-zuschauer, GE-C510-anfeuern, GE-C510-jubel, GE-C510-buhen

\chapter{Auf der Tribüne}
\label{ch:ge-auf-der-tribune}
\hlchaptermodality{510}

\begin{chapteropening}
\textbf{By the end of this chapter:} \emph{I can talk about watching sport in a stadium.}

Complete the last lesson of chapter 510: I can talk about watching sport in a stadium.
\end{chapteropening}

\section[die Tribüne]{die Tribüne — stand, grandstand}

\label{lesson:GE-C510-tribuene}



\begin{quote}
Before the new one: say the German for medal, then the German for award ceremony.
\end{quote}

\subsection*{You'll want to know: die Tribüne}
\textbf{die Tribüne} — \textquotedblleft{}stand, grandstand\textquotedblright{}.

\begin{quote}
Wir saßen auf der Tribüne direkt hinter dem Tor.
\end{quote}

\begin{quote}
\emph{We sat in the stand right behind the goal.}
\end{quote}

Plural: \emph{die Tribünen}.

\begin{cousinweb}
\emph{Tribüne} came from French \emph{tribune}, from Italian \emph{tribuna} (EtymDB), which is usually traced to Latin \emph{tribūnal}, the raised platform of the magistrates (Lewis \& Short).
\end{cousinweb}

\begin{grammarlens}[title={What you've built}]
One more word for telling, explaining and describing at B1.
\end{grammarlens}

\subsection*{Your turn}
\begin{itemize}
\raggedright
  \item \emph{Say it:} \emph{die Tribüne}
  \item \emph{Say it:} \emph{die Tribüne}, once more
  \item \emph{Say it:} a sentence of your own with \emph{die Tribüne}
  \item \emph{Recall:} say the German for medal, then the German for award ceremony, then say \emph{die Tribüne} again
\end{itemize}

\begin{tcolorbox}[breakable,colback=teal!4,colframe=teal!35!black,title={Before you move on}]
What does die Tribüne mean? (\textquotedblleft{}Stand, grandstand\textquotedblright{}.) Say it once more.
\end{tcolorbox}
\section[der Zuschauer]{der Zuschauer — spectator; viewer}

\label{lesson:GE-C510-zuschauer}



\begin{quote}
Before the new one: say the German for award ceremony, then the German for stand.
\end{quote}

\subsection*{You'll want to know: der Zuschauer}
\textbf{der Zuschauer} — \textquotedblleft{}spectator; viewer\textquotedblright{}.

\begin{quote}
Im Stadion waren fünfzigtausend Zuschauer.
\end{quote}

\begin{quote}
\emph{There were fifty thousand spectators in the stadium.}
\end{quote}

Plural: \emph{die Zuschauer}. Also a TV viewer.

\begin{cousinweb}
\emph{Zuschauer} is built on \emph{zuschauen} \textquotedblleft{}to watch\textquotedblright{}, from \emph{schauen} \textquotedblleft{}to look\textquotedblright{}, inherited from Old High German \emph{scouwōn}, from Proto-Germanic \emph{skawwōną} (EtymDB).
\end{cousinweb}

\begin{grammarlens}[title={What you've built}]
One more word for telling, explaining and describing at B1.
\end{grammarlens}

\subsection*{Your turn}
\begin{itemize}
\raggedright
  \item \emph{Say it:} \emph{der Zuschauer}
  \item \emph{Say it:} \emph{der Zuschauer}, once more
  \item \emph{Say it:} a sentence of your own with \emph{der Zuschauer}
  \item \emph{Recall:} say the German for award ceremony, then the German for stand, then say \emph{der Zuschauer} again
\end{itemize}

\begin{tcolorbox}[breakable,colback=teal!4,colframe=teal!35!black,title={Before you move on}]
What does der Zuschauer mean? (\textquotedblleft{}Spectator; viewer\textquotedblright{}.) Say it once more.
\end{tcolorbox}
\section[anfeuern]{anfeuern — to cheer on}

\label{lesson:GE-C510-anfeuern}



\begin{quote}
Before the new one: say the German for stand, then the German for spectator.
\end{quote}

\subsection*{You'll want to know: anfeuern}
\textbf{anfeuern} — \textquotedblleft{}to cheer on\textquotedblright{}.

\begin{quote}
Die Fans feuern ihre Mannschaft laut an.
\end{quote}

\begin{quote}
\emph{The fans loudly cheer on their team.}
\end{quote}

Separable: \emph{wir feuern an}, \emph{hat angefeuert}.

\begin{cousinweb}
\emph{anfeuern} is built on \emph{an} \textquotedblleft{}on\textquotedblright{} and \emph{feuern} \textquotedblleft{}to fire\textquotedblright{}, from \emph{Feuer} \textquotedblleft{}fire\textquotedblright{}, inherited from Old High German \emph{fiur} (EtymDB): to fire someone up.
\end{cousinweb}

\begin{grammarlens}[title={What you've built}]
One more word for telling, explaining and describing at B1.
\end{grammarlens}

\subsection*{Your turn}
\begin{itemize}
\raggedright
  \item \emph{Say it:} \emph{anfeuern}
  \item \emph{Say it:} \emph{anfeuern}, once more
  \item \emph{Say it:} a sentence of your own with \emph{anfeuern}
  \item \emph{Recall:} say the German for stand, then the German for spectator, then say \emph{anfeuern} again
\end{itemize}

\begin{tcolorbox}[breakable,colback=teal!4,colframe=teal!35!black,title={Before you move on}]
What does anfeuern mean? (\textquotedblleft{}To cheer on\textquotedblright{}.) Say it once more.
\end{tcolorbox}
\section[der Jubel]{der Jubel — cheering, jubilation}

\label{lesson:GE-C510-jubel}



\begin{quote}
Before the new one: say the German for spectator, then the German for to cheer on.
\end{quote}

\subsection*{You'll want to know: der Jubel}
\textbf{der Jubel} — \textquotedblleft{}cheering, jubilation\textquotedblright{}.

\begin{quote}
Nach dem Tor war der Jubel groß.
\end{quote}

\begin{quote}
\emph{After the goal there was huge cheering.}
\end{quote}

No plural. The verb: \emph{jubeln} — to cheer.

\begin{cousinweb}
\emph{Jubel} comes from Middle High German \emph{jūbel}, from Latin \emph{iūbilus} (EtymDB); Lewis \& Short gives \emph{iūbilum} \textquotedblleft{}a wild cry, a shout\textquotedblright{}.
\end{cousinweb}

\begin{grammarlens}[title={What you've built}]
One more word for telling, explaining and describing at B1.
\end{grammarlens}

\subsection*{Your turn}
\begin{itemize}
\raggedright
  \item \emph{Say it:} \emph{der Jubel}
  \item \emph{Say it:} \emph{der Jubel}, once more
  \item \emph{Say it:} a sentence of your own with \emph{der Jubel}
  \item \emph{Recall:} say the German for spectator, then the German for to cheer on, then say \emph{der Jubel} again
\end{itemize}

\begin{tcolorbox}[breakable,colback=teal!4,colframe=teal!35!black,title={Before you move on}]
What does der Jubel mean? (\textquotedblleft{}Cheering, jubilation\textquotedblright{}.) Say it once more.
\end{tcolorbox}
\section[buhen]{buhen — to boo}

\label{lesson:GE-C510-buhen}



\begin{quote}
Before the new one: say the German for to cheer on, then the German for cheering.
\end{quote}

\subsection*{You'll want to know: buhen}
\textbf{buhen} — \textquotedblleft{}to boo\textquotedblright{}.

\begin{quote}
Nach dem schlechten Spiel buhten die Zuschauer.
\end{quote}

\begin{quote}
\emph{After the poor game the spectators booed.}
\end{quote}

Regular: \emph{hat gebuht}. From the cry \emph{buh!}

\begin{cousinweb}
\emph{buhen} is built on the cry \emph{buh!}; EtymDB lists the verb without a source, and it is usually explained as imitative.
\end{cousinweb}

\begin{grammarlens}[title={What you've built}]
One more word for telling, explaining and describing at B1.
\end{grammarlens}

\subsection*{Your turn}
\begin{itemize}
\raggedright
  \item \emph{Say it:} \emph{buhen}
  \item \emph{Say it:} \emph{buhen}, once more
  \item \emph{Say it:} a sentence of your own with \emph{buhen}
  \item \emph{Recall:} say the German for to cheer on, then the German for cheering, then say \emph{buhen} again
\end{itemize}

\begin{tcolorbox}[breakable,colback=teal!4,colframe=teal!35!black,title={Before you move on}]
What does buhen mean? (\textquotedblleft{}To boo\textquotedblright{}.) Say it once more.
\end{tcolorbox}
